\begin{lemma}[The tail of a minimum hierarchy counterexample has at least two additions]
Fix a finite loopless multigraph with its individual edge identities, a
palette of $k$ colors, and a distinguished edge $e_0$.  Let
$(T,\varphi_n)$ be a positive-rung extended Tashkinov tree constructed
from an arbitrary certified input coloring, let $\mathcal B$ be a
hierarchy base, and let $C$ be a hierarchy counterexample over that base.
Write $H=C.\mathit{hierarchy}$, let $q$ be its divider count, and put
\[
 A=\noderef{HierarchyStarredStage}(H,q).
\]
Assume that $\varphi_n$ is proper and leaves exactly $e_0$ uncolored.

Assume, uniformly for every positive-rung maximum ETT on the same graph,
palette, and distinguished edge, with arbitrary certified input and
generating colorings and with rung number smaller than that of $T$, all
six lower-rung assertions of the technical theorem: elementarity,
terminal interchangeability, the two PE stability and stored-pair
assertions, the terminal PE component assertion, the defective-endpoint
order assertion, and the mod-coloring assertion.  Assume also that every
same-rung maximum ETT with a good hierarchy having fewer than $q$
dividers is elementary.

Assume that $C$ is lexicographically minimal among hierarchy
counterexamples with exactly $q$ dividers, allowing the input coloring,
generating coloring, ETT, base, and hierarchy to vary.  Thus no such
$C'$ has
\[
 \noderef{HierarchyCounterexampleLess}
   (\noderef{HierarchyCounterexampleMeasure}(C'),
    \noderef{HierarchyCounterexampleMeasure}(C)).
\]
Let $p_H=\noderef{HierarchyPathNumber}(H)$ and let
\[
 T(y_{p_H-1})=
 T\big|_{\min\{|A|+p_H-1,|T|\}}
\]
be the exact ordered prefix supplied by \noderef{TreeSequencePrefix}.
Assume (6.6), namely that this prefix is elementary under $\varphi_n$,
and assume the large-branch estimate
\[
 2n+11\leq |\overline\varphi_n(T_n)|,
\]
where $n$ is the rung number and $T_n$ is the terminal ladder tree.
Then
\[
                         2\leq |T|-|A|.
\]
Here $|T|-|A|$ is the number of exact edge/vertex additions above the
last starred stage.  It is not the path number $p_H$.
\end{lemma}
\begin{proof}
All tree inclusions in this proof are inclusions of ordered
edge-valued tree-sequences.  In particular, every retained edge is the
same element of the edge type, even when another parallel edge has the
same endpoints.

Put $\ell=|T|-|A|$.  We first record the ordered inclusion
\[
                              A\preccurlyeq T.                 \tag{1}
\]
If $q>0$, the positive-index clause of
\noderef{HierarchyStarredStage} identifies $A$ with the penultimate
hierarchy stage $H.\mathit{stage}(q)$, while the terminal-stage field of
\noderef{TashkinovHierarchy} identifies $T$ with
$H.\mathit{stage}(q+1)$.  Its `nested' field gives (1), and its `strict'
field gives $|A|<|T|$.  If $q=0$, the zero-index clause of
\noderef{HierarchyStarredStage} identifies $A$ with the core in
\noderef{HierarchyBase}; the `coreInsideFirst' and terminal-stage fields
of \noderef{TashkinovHierarchy} give (1).

We next exclude $\ell=0$.  There is nothing to prove when $q>0$, by the
strictness just noted.  Suppose $q=0$ and $\ell=0$.  Equality of lengths
in (1), followed by the vertex and edge clauses of
\noderef{TreeSequenceInitialSegment}, gives $A=T$ as complete ordered
tree-sequences.  Moreover, the candidate interval in
\noderef{HierarchyPathNumber} is then empty, so its fallback gives
$p_H=0$.  Truncated subtraction gives $p_H-1=0$, and hence the prefix in
the hypothesis has length
\[
 \min\{|A|+0,|T|\}=|T|.
\]
The full-length case of \noderef{TreeSequencePrefix}, together with
$A=T$, therefore identifies that prefix with both $A$ and $T$.  The
prefix-elementarity hypothesis makes $A$ elementary.  Applying
\noderef{HierarchyCounterexampleCoreProper} to this exact core says
$T\ne A$, contrary to the equality above.  Equivalently, the
`notElementary' field of \noderef{HierarchyCounterexample} contradicts
the same prefix elementarity.  Thus
\[
                              1\leq\ell.                       \tag{2}
\]

Assume for a contradiction that $\ell<2$.  By (2), $\ell=1$.  The
candidate interval in \noderef{HierarchyPathNumber} is now the singleton
$\{1\}$.  If its filtered candidate set is nonempty, its minimum is one;
if it is empty, the fallback tail length is one.  Thus
\[
                              p_H=1.                           \tag{3}
\]
Combining (1), $|T|=|A|+1$, and the defining vertex and edge equalities
of \noderef{TreeSequenceInitialSegment}, write the exact last addition as
\[
                  T=A\cup\{e_1,y_1\},                         \tag{4}
\]
where $e_1$ is the entry of $T$ at edge index $|A|$ and $y_1$ is the
new vertex entry at vertex index $|A|+1$.  No endpoint pair has been
substituted for $e_1$.  By (3), the prefix in (6.6) has length $|A|$;
the exact-prefix equality in \noderef{TreeSequencePrefix} and (1) make it
equal to $A$.  Consequently
\[
          V(A)\text{ is elementary under }\varphi_n.          \tag{5}
\]

We now derive the paper's common missing color rather than assuming it.
Negating the pairwise-disjointness statement in
\noderef{ColoredVertexSetElementary}, the `notElementary' field of $C$
supplies distinct vertices $u,v\in V(T)$ and a color $\alpha$ which is
missing at both.  By \noderef{TreeSequenceVertexSet} and the exact
one-vertex description (4),
\[
                         V(T)=V(A)\cup\{y_1\}.                 \tag{6}
\]
Both $u$ and $v$ cannot lie in $V(A)$, by (5), and they cannot both be
$y_1$, because they are distinct.  After interchanging their names,
$u\in V(A)$ and $v=y_1$.  Expanding \noderef{TreeMissingColors} as the
union of the sets from \noderef{MissingColors}, we obtain
\[
 \alpha\in\overline\varphi_n(A)
       \cap\overline\varphi_n(y_1).                           \tag{7}
\]

We split according to the divider count.

\emph{Suppose first that $q=0$.}
Let $n$ denote the last positive series index.  Positivity in
\noderef{HierarchyCounterexample} excludes the initial constructor when
we eliminate the certificate at $n$ from \noderef{TashkinovSeries}.
Thus the final certificate in \noderef{TashkinovSeriesStep} is exactly
one of SE, RE, and PE.

In either non-PE case, the `nonparallel' field of
\noderef{HierarchyBase} and the zero-index definition of $A$ give
\[
                              A=T_n.                           \tag{8}
\]
The SE and RE constructors first adjoin their individual connecting edge
$f_n$ to $T_n$ in the sense of \noderef{TreeSequenceAdjoinEdge} and only
then take a \noderef{TashkinovClosure}.  The exact first-edge conclusion
of \noderef{HierarchyNonparallelFirstEdge}, specialized to the
zero-divider hierarchy whose first stage is $T$, identifies the edge at
index $|T_n|$ with $f_n$.  Since (4) gives the same entry as $e_1$, and
the adjoin certificate fixes the new endpoint, entrywise equality gives
\[
                       f_n=e_1,\qquad u_n=y_1.                 \tag{9}
\]
This is the stored-seed exception: $e_1$ is not treated as an ordinary
Tashkinov augmentation in these two cases.

Assume that the final kind is SE.  Its
\noderef{MaximumDefectiveChoice} supplies the selected coloring and the
exact outside endpoint $u_n$.  The terminal `currentColoring' equation,
together with the generating-coloring field of the ETT, identifies that
selected coloring with $\varphi_n$.  Reflexivity satisfies both clauses
of \noderef{StableColoring}.  Hence the SE `seriesAlternative', applied
to $\varphi_n$, says that the missing-color set at $u_n$ is disjoint from
$\overline\varphi_n(T_n)$.  Equations (7)--(9) put $\alpha$ in both
sets, a contradiction.  Therefore the terminal kind is not SE.

Assume next that the final kind is RE.  Let
$(\gamma_n,\delta_n)$ and the exact cycle $O$ be the data in its
\noderef{RevisitingExtensionCertificate}.  We need the missingness of
$\delta_n$ on $T_n$, so we verify it from the stored certificate.  The
certificate chooses an earlier positive PE index $h$, records the exact
equations for its $\gamma_h,\delta_h$, and says every index strictly
between $h$ and $n$ is RE.  At that PE index,
\noderef{ParallelExtensionColor} makes $\gamma_h$ missing at the support
under the selected coloring.  The exact defective edge supplied by
\noderef{MaximumDefectiveChoice} is incident there and has color
$\delta_h$; the second \noderef{StableColoring} relation preserves that
edge and the missingness of $\gamma_h$.  The PE
\noderef{TwoColorComponentSwap} through the support therefore removes
the incident $\delta_h$ edge and replaces it by a $\gamma_h$ edge, so
$\delta_h$ is missing at the support under the PE output coloring.  Each
intervening RE constructor has `currentColoring' equal to its preceding
coloring, so induction through these exact equations preserves that
missing color at that same vertex.  The ladder trees are nested ordered
sequences, hence the vertex remains in $T_n$.  The equations
$\delta_h=\delta_n$ in the revisiting certificate give
\[
                       \delta_n\in\overline\varphi_n(T_n).    \tag{10}
\]
The terminal RE `previousClosed' field is a
\noderef{ColoredVertexSetClosed} assertion for $T_n$ under the previous
coloring, and its `currentColoring' field identifies that coloring with
$\varphi_n$.  Thus (7), (8), and (10) make both $\alpha$ and $\delta_n$
closed at $T_n$.

If $\alpha=\delta_n$, then $O$ already contradicts (7): by (9) its exact
edge $e_1$ lies on the $(\gamma_n,\delta_n)$-cycle, so the incidence
cycle has two incident colored edge-nodes at $y_1$, whereas
$\delta_n=\alpha$ is missing there.  Hence assume
$\alpha\ne\delta_n$.  Let $Q$ be the full
$(\delta_n,\alpha)$-component through $y_1$.  Closedness of both colors
at $T_n$ implies that $Q$ has no vertex in $V(T_n)$: a first entry into
$T_n$ would use an exact boundary edge colored $\alpha$ or $\delta_n$,
contrary to \noderef{ColoredVertexSetClosed}.  Since $y_1$ misses
$\alpha$, the component is a possibly singleton
\noderef{TwoColorPathComponent}; singleton components are retained here.

Let
\[
 \sigma_n=\noderef{TwoColorComponentSwap}
       (G,\varphi_n,\delta_n,\alpha,\operatorname{inl}y_1).
\]
The component $Q$ is disjoint from the component through the first
$\delta_n$-missing vertex of $T_n$.  Together with (7), (8), (10), and
$p_H=1$, this is exactly the zero-divider Case 2 certificate of
\noderef{StableChange}.  Its conclusion gives weak stability on $A$;
unfolding \noderef{HierarchyWeaklyStable} in the non-PE case makes this
the precise
$(T_n,D_n,\varphi_n)$-\noderef{StableColoring} relation.  It also gives
properness and the same uniquely uncolored edge.  Applying
\noderef{Td} to this stable coloring and the stated six lower-rung
package gives
\[
                       \noderef{ModColoring}(T,\sigma_n).      \tag{11}
\]

The exact one-edge seed in (9) is $T$ itself, so the reflexive
augmentation sequence after that seed satisfies the RE alternative in
(11).  The universal clause of \noderef{ModColoring} therefore supplies
an actual \noderef{CorrespondingExtendedTashkinovTree} on the same
ordered $T$.  Its extension-kind and connecting-edge-set agreement say
that its terminal step is again RE and its exact connecting set is
$\{e_1\}$.  Eliminate that replacement RE constructor.  Its own
revisiting certificate supplies a two-color cycle through its connecting
edge.  Equality of raw stored pairs says that the unordered pair of its
cycle colors is exactly $\{\gamma_n,\delta_n\}$; no unjustified equality
of the two ordered internal witnesses is used.  The swap was outside
$T_n$, and $y_1$ missed $\alpha$ before it, so $y_1$ misses
$\delta_n$ under $\sigma_n$.  The replacement cycle contains the exact
edge $e_1$ and therefore the endpoint $y_1$, but a
\noderef{TwoColorCycle} has degree two at every node and in particular
uses both of its two colors at a graph-vertex node.  This is impossible
when one member $\delta_n$ of its color pair is missing at $y_1$.
Thus the terminal kind is not RE.

It remains in the zero-divider case to consider PE.  The parallel field
of \noderef{HierarchyBase} identifies
\[
                 A=T_n\vee R_n                               \tag{12}
\]
as the exact ordered union with the stored closure $R_n$.  In a PE there
is no exceptional SE/RE seed: the final
\noderef{TashkinovClosure} is obtained from $T_n$ by a
\noderef{TashkinovAugmentationSequence}.  Its one-at-a-time length
increments pass through the prefix lengths $|A|$ and $|T|=|A|+1$, so
the intervening \noderef{TashkinovAugmentation} is exactly the displayed
edge $e_1$.  Consequently $e_1$ has some color $\beta$ missing on $A$.
Because $e_1$ is incident with $y_1$ while $\alpha$ is missing there,
$\alpha\ne\beta$.

The `parallelBase' field of \noderef{GoodHierarchy} supplies at least
four colors in
\[
 (\overline\varphi_n(T_n)\cap\overline\varphi_n(R_n))
       -\bigcup_{\eta\in D_{n,0}}\Gamma^0_\eta.
\]
Choose $\theta$ from this set distinct from $\alpha$ and $\beta$.
The elementarity premise needed below is obtained without manufacturing
a zero-rung ETT.  If $n=1$, it is the first conclusion of
\noderef{TechnicalTheoremBase}.  If $n>1$, restrict the certified series
through index $n-1$ and display its exact last closure $T_n$; the
quantified fields of \noderef{MaximumProperty} restrict with the series, so lower
assertion (i) applies.  Thus the hypothesis of
\noderef{ParallelCommonMissingClosed} holds, and that lemma makes
$\theta$ closed on the union (12).  In the notation of
\noderef{TreeColorSetClosed},
\[
       \theta\in\overline\varphi_n(A),\qquad
       \partial_{\varphi_n,\theta}(A)=\varnothing.            \tag{13}
\]

Apply conclusion (iii) of \noderef{interchange} with first color
$\theta$ and second color $\alpha$.  Its combined-stability hypothesis
is reflexive under $\varphi_n$, and (13) supplies its missing and closed
premises.  Hence $\theta$ and $\alpha$ are $A$-interchangeable.  Choose
$x_\theta,x_\alpha\in V(A)$ missing the indicated colors.  If they
coincide, their two-color component is the singleton vertex component.
If they are distinct, properness and
\noderef{TwoColorIncidenceDegreeAtMostTwo}, together with
\noderef{TwoColorComponentSpanningPath}, make them the two ends of their
component.  In either case this is a path component meeting $A$.  If the
component through $y_1$ also met $A$, interchangeability would identify
it with this component; in the coincident case that is impossible because
the component is singleton, and in the distinct case it would add the
third end $y_1$.  Thus the possibly singleton
$(\theta,\alpha)$-component $Q$ through $y_1$ is disjoint from $A$.
This uses exactly \noderef{TreeInterchangeable}, whose definition retains
singleton components.

Swap $Q$ and again write the result as $\sigma_n$, now with ordered
colors $(\theta,\alpha)$.  The colors are both missing on the prefix
$A$, $\theta$ is missing on $T_n$, and $Q$ is disjoint from the
component through the first $\theta$-missing vertex.  Thus the
zero-divider Case 2 certificate of \noderef{StableChange} applies.  Its
weak-stability conclusion unfolds via \noderef{HierarchyWeaklyStable} in
the PE case to combined stability on $T_n$ and $R_n$.  The edge and
vertex clauses of \noderef{StableColoring} fix every exact edge incident
with their union and preserve all missing sets there.  In particular,
$\theta$ remains missing and closed on $A$, $\beta$ remains missing on
$A$, and $e_1$ retains color $\beta$; meanwhile $y_1$ now misses
$\theta$.

Apply conclusion (iii) of \noderef{interchange} under $\sigma_n$ to
$\theta,\beta$.  It says these colors are $A$-interchangeable.  Choose
witnesses in $A$ missing $\theta$ and $\beta$.  If they coincide, their
component is the singleton at that vertex and is distinct from the
component containing $y_1$ and the exact edge $e_1$.  If they are
distinct, they are two ends of a spanning path, while $y_1$ is a third
end of the component containing $e_1$ if the components were equal.
Thus at least two path components meet $A$, contradicting
\noderef{TreeInterchangeable}.  All three terminal kinds are impossible,
so $q=0$ is impossible under the assumption $\ell=1$.

\emph{Suppose now that $q>0$.}
Then $A=H.\mathit{stage}(q)$.  The exact last edge $e_1$ is an ordinary
Tashkinov augmentation, not the stored-seed exception.  Indeed,
\noderef{HierarchyNonparallelFirstEdge} puts the individual SE/RE seed
edge, when present, in the first positive hierarchy stage, hence already
in $A$.  The finite relation in
\noderef{TashkinovAugmentationSequence} for the terminal
\noderef{TashkinovClosure} increases the length by one at each
\noderef{TashkinovAugmentation}.  Both $A$ and $T$ are ordered prefixes
of its final tree and their lengths differ by one, so the sequence step
between those two lengths is exactly $(e_1,y_1)$.  Let
\[
                         \varphi_n(e_1)=\gamma.                \tag{14}
\]
The missing-color witness in \noderef{TashkinovAugmentation} gives
$\gamma\in\overline\varphi_n(A)$.  The terminal `minusClosed' field of
\noderef{GoodHierarchy}, expanded using
\noderef{TreeColorSetMinusClosed}, therefore forces
\[
 \gamma\in
   \bigcup_{\eta\in D_{n,q}}\Gamma^{q-1}_\eta;               \tag{15}
\]
otherwise the boundary edge $e_1$ of the missing color $\gamma$ would
be forbidden.

Let $A^-$ be the preceding starred stage.  Strict hierarchy growth gives
a vertex $x\in V(A)-V(A^-)$.  The initial $k$-triple inside the certified
\noderef{TashkinovSeries} has $k\geq\Delta(G)+1$ by
\noderef{KTriple}.  By \noderef{Degree} and \noderef{MaximumDegree},
at most $\Delta(G)$ individual edges are incident with $x$; and
\noderef{IsProperEdgeColoring} uses at most one of them for each color.
Hence
some color $\theta$ is missing at $x$.  Elementarity (5) implies that
$\theta$ is not missing at any vertex of $A^-$, since $x$ is outside
that set.  Thus $\theta$ is newly missing above the last divider.
Applying \noderef{GoodHierarchyNewMissingClosed} gives
\[
 \theta\in\overline\varphi_n(A),\qquad
 \partial_{\varphi_n,\theta}(A)=\varnothing.                 \tag{16}
\]
By \noderef{HierarchyStoredColorsAntitone}, each
$\eta\in D_{n,q}$ also belongs to the preceding stored set.  The
avoidance clause of \noderef{GoodHierarchy} then puts every color in the
union (15) in the preceding \noderef{HierarchyStageMissingColors}; its
definition supplies a missing-color witness in $A^-$.  Since $\theta$
has no such witness, $\theta\ne\gamma$.

If $\theta=\alpha$, apply \noderef{ChangeLemma} directly to the colors
$\theta,\gamma$.  By (3) its prefix is $A$; (5), (14), and (16) give
all missing and closed premises, while the ordered-unused implication is
vacuous because both colors are missing on $A$.  The conclusion makes
$\theta,\gamma$ $A$-interchangeable.  Choose witnesses in $A$ missing
$\theta$ and $\gamma$.  If they coincide, their two-color component is
the isolated vertex component and is distinct from the component
containing $e_1$ and $y_1$.  If they are distinct, they are two ends of
a spanning path, and equality with the component through $y_1$ would add
a third end.  Either alternative contradicts
\noderef{TreeInterchangeable}.  We may therefore assume
\[
                         \theta\ne\alpha.                     \tag{17}
\]

Apply \noderef{step1} at the last positive hierarchy stage with first
color $\theta$ and second color $\alpha$.  Its hypotheses are exactly
the stated lower six-pack, fewer-divider induction, large estimate,
maximum property and positivity from $C$, and (16)--(17).  It follows
that $\theta$ and $\alpha$ are $A$-interchangeable.  The same endpoint
argument used in the PE case shows that the full
$(\theta,\alpha)$-component $Q$ through $y_1$ has no vertex in $A$ and
is a possibly singleton \noderef{TwoColorPathComponent}.

Set
\[
 \sigma_n=\noderef{TwoColorComponentSwap}
       (G,\varphi_n,\theta,\alpha,\operatorname{inl}y_1).     \tag{18}
\]
For clarity, the general outside-hierarchy interface also applies to the
prefix at this point.  Use \noderef{GoodHierarchyPrefixRepackage} at
stage $A$ to obtain an actual maximum ETT on the same exact ordered
prefix, with the retained good hierarchy and exact connecting-edge
identities.  Since $Q$ has no vertex in $A$, (18) is an
\noderef{OutsideKempeRecoloring} of that prefix.  Applying
\noderef{HierarchyLEMStable} gives properness, unique uncoloredness,
stability and invariance on $A$, and an actual corresponding maximum ETT
with the same retained good hierarchy.  This invocation treats an
isolated $Q$ as a path component and changes colors only on its individual
edge-nodes.

We also need the stronger terminal data recorded specifically for this
minimum counterexample.  In \noderef{StableChange}, take its positive-
divider Case 1 with first color $\theta$, second color $\alpha$, and root
$y_1$.  Its hypotheses are (3), (5), (16)--(18), the path-component
certificate just proved, and disjointness from the component through the
first $\theta$-missing vertex.  Its conclusion gives
\[
 \begin{gathered}
  \sigma_n\text{ is weakly stable on }A,
  \qquad \sigma_n\text{ is a mod-coloring of the original ETT},\tag{19}\\
  A\text{ carries the transported maximum prefix hierarchy, and is}
  \left(\bigcup_{\eta\in D_{n,q}}\Gamma^{q-1}_\eta\right)^-
  \text{-closed under }\sigma_n.                             \tag{20}
 \end{gathered}
\]
The retained $D$-sets and every paper Gamma pair agree exactly.

We spell out why (19)--(20) also reconstruct the full displayed $T$.
The edge $e_1$ is not in $Q$, because it is incident with $A$ whereas
$Q$ has no vertex there; furthermore its color $\gamma$ is distinct from
$\theta$ by (16) and from $\alpha$ by (7) and its incidence with $y_1$.
Thus (14) is unchanged.  The weak-stability and invariance clauses in
\noderef{HierarchyWeaklyStable} preserve every earlier augmentation
from the terminal ladder tree through $A$, including the exact SE/RE
seed edge if that exception occurs, and preserve each old missing-color
witness.  Append the same exact augmentation $(e_1,y_1)$, whose color
$\gamma$ is still missing on $A$.  This gives the precise PE or exact
SE/RE eligibility alternative for $T$ in the universal clause of the
\noderef{ModColoring} assertion (19).  Hence there is a corresponding
ETT on the same ordered $T$.  Applying
\noderef{MaximumPropertyTransport} makes it maximum.

Use the transported base on $A$, keep every exact old stage and divider,
and append the unchanged terminal stage $T$.  Pair cardinalities and
disjointness are literal copies.  Avoidance and unique refresh through
the last divider use only missing colors on $A$ and the earlier starred
stages, all preserved by weak stability, together with the same stored
sets and Gamma pairs.  At the final Gamma index, the only displayed
segment edge above $A$ is $e_1$; it is fixed by (18), so the final
segment-color set is unchanged.  The last minus-closed field is exactly
(20), and all earlier minus-closed fields belong to the transported
prefix.  Thus these data form a good hierarchy with the same
divider count and the full equality relation expressed by
\noderef{SameGoodHierarchy}.  Under the swap, $y_1$ now misses $\theta$,
while stability on $A$ preserves the vertex of $A$ missing $\theta$.
The full replacement tree is therefore not elementary.  We have obtained
a hierarchy counterexample $C_\sigma$ on the same ordered $T$.

Its measure equals that of $C$: the tree and the complete hierarchy
stages, including every exact edge entry, are unchanged.  Consequently
the stated minimality of $C$ is also minimality for $C_\sigma$ after
rewriting by this equality of
\noderef{HierarchyCounterexampleMeasure}.  The path number is still one,
the relevant prefix is still $A$, and stability preserves its
elementarity.  It also preserves the missing-color set of $T_n$, so the
large estimate remains valid.  Thus every ambient hypothesis of
\noderef{ChangeLemma} holds for $C_\sigma$.

Both $\theta$ and $\gamma$ are missing on $A$ under $\sigma_n$.
Moreover, (15), (16), and (20) show that $\theta$ lies outside the
reserved union and is closed on $A$ under $\sigma_n$, after unfolding
\noderef{TreeColorSetMinusClosed} and \noderef{TreeColorSetClosed}.
Apply case (iii) of \noderef{ChangeLemma} to $\theta,\gamma$; its
ordered-unused premise is again vacuous because both colors are missing
on $A$.  It says that $\theta$ and $\gamma$ are $A$-interchangeable under
$\sigma_n$.

This is the final contradiction.  The swap makes $\theta$ missing at
$y_1$, and the unchanged exact edge $e_1$ has color $\gamma$, so its
$(\theta,\gamma)$-path meets $A$.  Choose witnesses in $A$ missing
$\theta$ and $\gamma$.  If they coincide, their component is isolated
and cannot equal the component containing $e_1$.  If they are distinct,
they are two ends of a spanning path, so equality with the component
having the additional end $y_1$ violates
\noderef{TwoColorIncidenceDegreeAtMostTwo}.  Thus at least two path
components meet $A$, contradicting \noderef{TreeInterchangeable}.

Both $q=0$ and $q>0$ have contradicted the assumption $\ell=1$.  Together
with (2), this proves $2\leq\ell=|T|-|A|$, as required.
\end{proof}
